% !TEX root = ../../main.tex
\subsection{消息加密存储与哈希摘要机制}
\label{subsec:msg_crypto}

\subsubsection*{消息加密存储}

通信消息在写入数据库之前，经过 SM4-CBC 加密处理，数据库中不存储任何明文内容。消息处理的完整流水线如图~\ref{fig:msg_pipeline} 所示，加密与哈希计算并行进行，随后一并写入数据库。

% !TEX root = ../../../main.tex
% 消息处理流水线图
\begin{figure}[H]
\centering
\begin{tikzpicture}[
  node distance = 0.5cm and 1.0cm,
  pstep/.style = {rectangle, draw, rounded corners=3pt, minimum width=2.8cm, minimum height=0.72cm,
                  font=\small, fill=blue!8, draw=blue!50!black},
  pside/.style = {rectangle, draw, rounded corners=2pt, minimum width=3.0cm, minimum height=0.62cm,
                  font=\small, fill=orange!14, draw=orange!50!black},
  pout/.style  = {rectangle, draw, rounded corners=8pt, minimum width=2.6cm, minimum height=0.62cm,
                  font=\small, fill=green!14, draw=green!50!black},
  arr/.style   = {-Stealth, thick},
  lbl/.style   = {font=\small},
]

% 主流水线（横向）
\node[pstep]                        (recv)   {接收明文消息 $M$};
\node[pstep, right=1.0cm of recv]   (split)  {并行处理};
\node[pstep, right=1.0cm of split]  (store)  {写入 messages 表};
\node[pstep, right=1.0cm of store]  (merkle) {更新 Merkle Tree};
\node[pout,  right=1.0cm of merkle] (bcast)  {广播给房间成员};

% 并行分支（向上/向下）
\node[pside, above=1.0cm of split] (enc)  {SM4-CBC 加密 $\to$ 密文 $C$};
\node[pside, below=1.0cm of split] (hash) {SM3 摘要 $\to$ $h = \mathrm{SM3}(\cdots)$};

% 箭头
\draw[arr] (recv)   -- (split);
\draw[arr] (split.north) -- (enc.south);
\draw[arr] (split.south) -- (hash.north);
\draw[arr] (enc.east)  -- ++(0.4,0) |- (store.north);
\draw[arr] (hash.east) -- ++(0.4,0) |- (store.south);
\draw[arr] (store)  -- (merkle);
\draw[arr] (merkle) -- (bcast);

% 旁注
\node[lbl, above=0.1cm of enc]  {随机 IV，PKCS7 填充};
\node[lbl, below=0.1cm of hash] {orderID|role|ts|ciphertext};

\end{tikzpicture}
\caption{消息处理流水线（加密存储与哈希计算并行）}
\label{fig:msg_pipeline}
\end{figure}


加密流程如下：

\begin{enumerate}[label=\arabic*., leftmargin=2em]
  \item 系统从环境变量 \texttt{FOODSHIELD\_SM4\_KEY} 读取 SM4 密钥（128 位），若未配置则使用系统内置默认密钥；
  \item 使用 \texttt{secrets.token\_bytes(16)} 生成随机 IV（初始向量），每条消息使用独立 IV；
  \item 对明文字节串执行 PKCS\#7 填充，使数据长度对齐至 16 字节块边界；
  \item 执行 SM4-CBC 加密，得到密文$C$；
  \item 将$IV \| C$拼接后以十六进制字符串形式存入 \texttt{messages.content} 字段。
\end{enumerate}

存储格式可表示为：
\[
  \texttt{content} = \operatorname{hex}\!\left(\mathit{IV} \,\|\, \operatorname{SM4\text{-}CBC}_{K}(\mathit{IV},\, M)\right)
\]
其中$\mathit{IV}$占前 16 字节（32 位十六进制字符），后续为密文。解密时，系统从 \texttt{content} 字段读取十六进制字符串，提取前 16 字节作为$\mathit{IV}$，其余部分作为密文执行 SM4-CBC 解密，并去除 PKCS\#7 填充后还原明文。

每条消息使用独立随机 IV 的设计保证了相同明文在不同时间加密后产生不同密文，避免攻击者通过密文相等推断消息内容是否相同。

\subsubsection*{消息哈希摘要计算}

每条消息在写入数据库前，服务端在内存中对消息的明文字段进行联合哈希计算，生成 \texttt{message\_hash}，该哈希值作为 Merkle Tree 的叶子节点参与完整性验证。

消息哈希公式为：
\[
  h_i = \operatorname{SM3}(\mathit{orderID} \| \mathit{senderPID} \| \mathit{role} \| \mathit{content}_{\mathrm{plain}} \| \mathit{timestamp})
\]
其中各字段之间使用分隔符 \texttt{|} 拼接，防止字段边界歧义。$\mathit{content}_{\mathrm{plain}}$为消息明文，哈希计算在服务端内存中完成，计算完成后明文即被 SM4 加密替换，数据库中不保存明文摘要。

上述设计满足以下安全目标：
\begin{itemize}[leftmargin=2em]
  \item \textbf{内容绑定}：哈希值绑定了订单号、发送方身份、角色、明文内容和时间戳，任意字段被篡改都会导致重新计算的哈希与存储值不一致；
  \item \textbf{明文隔离}：哈希基于明文计算，但数据库中仅存储密文，攻击者即使获取数据库内容，也无法从 \texttt{content} 字段直接还原哈希输入；
  \item \textbf{完整性可验证}：后续验证时，服务端先对 \texttt{content} 字段执行 SM4 解密得到明文，再按相同公式重算哈希，与存储的 \texttt{message\_hash} 对比，实现单条消息层面的完整性验证。
\end{itemize}

\subsubsection*{管理员访问控制}

管理员查询消息记录时，消息查询接口仅返回\texttt{msg\_id}、\texttt{order\_id}、\texttt{sender\_pid}、\texttt{role}、\texttt{message\_hash}和\texttt{timestamp}字段，不返回\texttt{content}密文字段。Merkle 完整性验证只需\texttt{message\_hash}列表，管理员无需接触消息密文即可完成日志验证，实现审计与明文访问的权限隔离。
